% ==============================================================================
% SEZIONE 7: CONCLUSIONI E SINTESI METODOLOGICA
% ==============================================================================

\section{Conclusioni}

\begin{frame}{Sintesi Metodologica: Il Nuovo Framework}
    L'analisi del caso Enel dimostra che la valutazione di una utility moderna deve trascendere il dato finanziario isolato per abbracciare una metrica multidimensionale.
    \vspace{0.1cm}
    
    \note{Il caso Enel dimostra che la valutazione di una utility moderna richiede metriche integrate e multidimensionali.}

    \begin{block}{La Formula del Valore Integrato}
        La solidità dell'EBITDA (\textbf{€22,9 miliardi}) non è un dato esogeno, ma è il risultato sinergico di due fattori abilitanti:
        \begin{itemize}
            \item \textbf{Qualità degli Asset Regolati (RAB):} Fornisce flussi di cassa stabili e protetti.
            \item \textbf{Infrastruttura Digitale (Smart Metering):} Abbate il rischio di credito e fidelizza l'utente.
    \end{itemize}
    \end{block}
    \vspace{-0.1cm}
    \pause
    \note{La tenuta dell'EBITDA è legata alla qualità della RAB e alla capacità di digitalizzare la base clienti tramite smart metering.}

    \begin{alertblock}{Il Futuro del Settore}
        La transformación da utility tradizionale a gestore di piattaforme digitali indica che il valore risiede nella capacità di coniugare \textbf{rigore finanziario} e \textbf{centralità del dato tecnologico}.
    \end{alertblock}
    \note{Il futuro del settore risiede nella capacità di coniugare rigore finanziario e automazione digitale delle infrastrutture.}
\end{frame}


\begin{frame}{Linee Guida per l'Analista Finanziario}
    Per replicare questo framework analitico su altri player globali del settore energetico, si raccomandano tre direttrici tecniche fondamentali:
    \vspace{0.2cm}

    \begin{enumerate}
        \item \textbf{Analisi della Visibilità Regolatoria:} Valutare la stabilità del quadro normativo nazionale e il suo impatto sul \textbf{WACC (Costo Medio Ponderato del Capitale)} in relazione alla crescita della RAB.
        \pause
        \note{Per replicare l'analisi si raccomandano tre direttrici. Prima: valutare la stabilità regolatoria e l'impatto sul WACC.}
        
        \item \textbf{Isolamento dalla Volatilità:} Verificare il \textit{Commodity Insulation Rate}, misurando la quota di margine protetta da asset regolati rispetto alla generazione \textit{merchant} pura.
        \pause
        \note{Seconda: calcolare il commodity insulation rate per isolare il margine dai rischi della generazione merchant.}
        
        \item \textbf{Monetizzazione dei Servizi Digitali (VAS):} Misurare come l'adozione di tecnologie smart (2G, BESS) riduca il costo di acquisizione clienti (\textit{CAC}) e il costo del rischio (\textit{Bad Debt}).
    \end{enumerate}
    \note{Terza: verificare come le tecnologie smart e i sistemi BESS riducano il CAC e l'impatto del bad debt.}
\end{frame}
